% ---- Sources, method, and reproducibility (back matter) ----
% Honest statement of what is measured, what is derived, and how to reproduce it.
\chapter{Sources, Method, and Reproducibility}
\label{chap:sources}
\markboth{Sources, Method, and Reproducibility}{Sources, Method, and Reproducibility}

\noindent This handbook marks every number with a two-axis evidence label so a reader never has to guess the provenance or the epistemic status of a figure. Labels of the form \textsc{[provenance][epistemic-status]} combine one item from each axis.

\medskip
\noindent\textbf{Axis 1 --- Source provenance:} where the number comes from.

\begin{itemize}\tightlist
\item \textbf{\textsc{[1P]} First-party / primary source.} A vendor technical report, model card, or official documentation. These are the model's own stated parameters, benchmarks, and limits.
\item \textbf{\textsc{[2\ensuremath{^\circ}]} Secondary source.} A third-party measurement, survey, or synthesis (e.g. an independent benchmark or a community write-up). Requires independent verification.
\item \textbf{\textsc{[LAB]}} Measured on the author's own hardware. These are home-lab measurements, not production hardware, and are used to illustrate mechanisms and orders of magnitude.
\item \textbf{\textsc{[ILLUSTRATIVE]}} A worked-example or scenario input chosen to demonstrate a calculation, not a reported measurement. It is a number we assumed so the arithmetic would be concrete.
\end{itemize}

\medskip
\noindent\textbf{Axis 2 --- Epistemic status:} what the number claims to be.

\begin{itemize}\tightlist
\item \textbf{\textsc{[FACT]}} Stated by the source as a fact (e.g. a parameter count or a vendor benchmark).
\item \textbf{\textsc{[DERIVED]}} Computed by arithmetic from stated inputs. The derivation is our work; the inputs come from a source item.
\item \textbf{\textsc{[HYPOTHESIS]}} A reasoned expectation that has not yet been validated. Deliberately kept, because it distinguishes what we know from what we only expect.
\item \textbf{\textsc{[UNRESOLVED]}} An open item that could not be pinned to a first-party source at the time of writing. Such items are flagged in the text and \emph{must be re-checked against current vendor documentation before any procurement or deployment decision}; they are precisely the numbers most likely to have drifted.
\end{itemize}

\medskip
\noindent\textbf{Reading a two-axis label.} \textsc{[1P][FACT]} means a vendor-stated fact; \textsc{[1P][DERIVED]} means arithmetic we computed from a first-party input; \textsc{[LAB][FACT]} means a home-lab measurement; \textsc{[ILLUSTRATIVE][ASSUMPTION]} means a number we chose so the worked example would be concrete — it is not evidence of anything about the real workload. When only one axis is shown, the other is implied by the surrounding sentence (e.g. a bare \textsc{[DERIVED]} carries the status, with the provenance stated in prose).

\noindent\textbf{Prices, models, and benchmarks are a moving target.} The model names, parameter counts, context limits, prices, and throughput figures in this book are a dated snapshot. They vary with the provider, region, commit/on-demand/spot status, hardware SKU, precision and sparsity mode, quantization, batch and concurrency, software stack, GPU count and topology, kernel version, warmup procedure, and percentile method used. \textbf{Re-verify before acting.}

\noindent\textbf{Reproducibility.} The book's quantitative examples are worked from the stated inputs in the surrounding text; the home-lab measurements were made on the author's own hardware under the configuration described in the chapter. Where a number is flagged as unresolved, the workspace did not have a first-party source at hand, and we say so rather than present an estimate as fact.

\noindent\textbf{Home-lab scope.} The author's lab is a CPU-only workstation (Intel i5-3210M, four cores, 15\,GB RAM). Measurements reflect that environment and the reduced-precision / small-model regime; they are used to \emph{illustrate} mechanisms and orders of magnitude, not to characterize production hardware. Kernel-performance and throughput figures from public sources are to be treated as vendor-reported until re-measured on the target cluster. \textbf{No H100/H200 production-performance claim in this book should be interpreted as a home-lab measurement unless explicitly labeled as such.}

\noindent\textbf{Code and data availability.} The build scripts that produce this book (text $\to$ LaTeX $\to$ PDF, plus the figure generators) are in the accompanying repository. The figures are generated where possible (matplotlib and a hand-drawn diagram pipeline), so the arithmetic and layouts are inspectable rather than opaque artifacts.
